\documentclass[11pt]{article}
\usepackage[T1]{fontenc}
\usepackage{lmodern,amsmath,amsthm,microtype}
\usepackage[margin=1in]{geometry}
\usepackage{xcolor}
\definecolor{linkblue}{RGB}{28,66,116}
\usepackage[colorlinks=true,allcolors=linkblue,breaklinks=true]{hyperref}
\hypersetup{pdftitle={Subexponential Set Splitting from Exact SAT Counting},
pdfauthor={GPT-6 (OpenAI Codex)},
pdfsubject={A variable-efficient counting reduction and conditional subexponential algorithms},
pdfkeywords={set splitting, exact counting, satisfiability, binomial moments, parameterized algorithms}}
\usepackage{fancyhdr}
\pagestyle{fancy}
\fancyhf{}
\fancyhead[L]{\small Subexponential Set Splitting}
\fancyhead[R]{\small Research draft}
\fancyfoot[C]{\thepage}
\setlength{\headheight}{14pt}
\setlength{\emergencystretch}{2em}
\newtheorem{theorem}{Theorem}[section]
\newtheorem{lemma}[theorem]{Lemma}
\newtheorem{corollary}[theorem]{Corollary}
\theoremstyle{definition}
\newtheorem{definition}[theorem]{Definition}
\newcommand{\poly}{\operatorname{poly}}
\newcommand{\Count}{\mathsf{Count}}
\newcommand{\bits}{\{0,1\}}
\title{Subexponential Set Splitting\\from Exact SAT Counting}
\author{GPT-6\\\small OpenAI Codex}
\date{29 October 2026}

\begin{document}
\maketitle
\vspace{-1em}
\begin{abstract}
We reduce Set Splitting to exact counting for 3-CNF formulas.
A hypergraph with $n$ vertices and $m$ edges needs $m+1$ counting calls,
each with $O(n+m)$ variables, regardless of edge size.
The counts determine how many colorings split each possible number of edges.
Using the general counting bound stated in the October 2026 SAT manuscript,
this gives a deterministic polynomial-space algorithm with running time
\[
 |I|^{O(1)}+2^{O(k\log\log k/\sqrt{\log k})},
\]
where $k$ is the number of edges to split.
We state separately the weaker bound supported by its formalization scope
note. The reduction is proved here; the SAT counting result is assumed.
\end{abstract}

\section{Problem and assumptions}

Set Splitting asks whether a two-coloring can make at least $k$ sets
contain both colors. Sets may have any size.
Lokshtanov and Saurabh~\cite{LS09} give an equivalent instance with at most
$k$ vertices and $2k$ edges.
Earlier specialized algorithms take $O^*(1.7171^k)$ time with exponential
space or $O^*(1.8213^k)$ time with polynomial space~\cite{NRD14}.
Here $O^*$ hides polynomial factors.

Let $H=(V,(e_i)_{i=1}^m)$ be a hypergraph with $V=[n]$.
Edges have separate labels, so repeated edges count separately.
Write
\[
 s=n+m,\qquad J=\sum_i |e_i|.
\]
The vertices and edges are listed explicitly; $L$ is the input bit length.
Empty and singleton edges are never split.
For this paper, both count as monochromatic.

For $\alpha\in\{1,2\}$, define
\[
 f_\alpha(v)=\frac{v(\log\log v)^\alpha}{\sqrt{\log v}}.
\]
Small arguments are handled directly.
We use the following external assumption.

\begin{definition}[SAT counting assumption]\label{def:count}
$\Count_\alpha$ is a uniform deterministic algorithm that counts the
satisfying assignments of a 3-CNF with $v$ declared variables and encoding
length $\ell$ in time
\[
 (\ell+1)^{O(1)}2^{O(f_\alpha(v))}
\]
and polynomial space.
\end{definition}

The revised SAT manuscript states $\Count_1$~\cite{SAT26}.
Its formalization scope note documents $\Count_2$~\cite{Scope26}.
We keep these rates separate and do not reprove either result.
All formulas below have width at most three.
Unused variables remain declared and contribute their usual factor of two.

\section{Counting selected monochromatic edges}

For a coloring $x\in\bits^n$, let $u(x)$ be the number of monochromatic
edges. Define
\[
 a_j=|\{x:u(x)=j\}|,\qquad
 P_H(z)=\sum_{j=0}^m a_j z^{m-j}.
\]
Thus the coefficient of $z^b$ counts colorings that split exactly $b$ edges.
This is a form of the hypergraph coboundary polynomial~\cite{White11}.

We first compute
\begin{equation}\label{eq:moments}
 F_t=\sum_{x\in\bits^n}\binom{u(x)}t,\qquad 0\le t\le m.
\end{equation}
These counts select $t$ monochromatic edges, then sum over all colorings.

\begin{lemma}[Recovering the distribution]\label{lem:inverse}
The counts in \eqref{eq:moments} give
\begin{equation}\label{eq:inverse}
 a_j=\sum_{t=j}^m(-1)^{t-j}\binom tj F_t.
\end{equation}
This calculation takes polynomial time and space.
\end{lemma}
\begin{proof}
The binomial theorem gives
\[
 \sum_{t=0}^m F_t w^t
 =\sum_x(1+w)^{u(x)}
 =\sum_{j=0}^m a_j(1+w)^j.
\]
Replace $w$ by $z-1$ and compare coefficients.
Each $F_t\le2^{n+m}$.
Each term in \eqref{eq:inverse} has magnitude at most $2^{n+2m}$, so
partial sums need only $O(n+m+\log(m+1))$ bits.
There are $O(m^2)$ terms.
\end{proof}

\begin{lemma}[Counting selected bits]\label{lem:adder}
The condition that exactly $t$ of $m$ bits equal one has a width-three
CNF encoding with $O(m+1)$ extra variables and clauses.
Each valid input assignment has exactly one extension to the extra variables.
\end{lemma}
\begin{proof}
Sum the bits using a balanced tree of binary adders, then fix the sum to $t$.
At level $h$, there are $O(m/2^h)$ additions of $O(h)$-bit numbers.
The gate count is
\[
 O\!\left(m\sum_{h\ge1}\frac{h}{2^h}\right)=O(m).
\]
Padding with zeros handles other input lengths.
Encode each gate by its full logical equivalence.
For example, $g\leftrightarrow(a\wedge b)$ needs
\[
 (\neg g\vee a)\wedge(\neg g\vee b)\wedge(g\vee\neg a\vee\neg b).
\]
OR, NOT, and XOR also have constant-size width-three encodings.
The inputs determine every gate value, giving the required unique extension.
The case $m=0$ takes constant space.
\end{proof}

\begin{theorem}[SAT reduction]\label{thm:encoding}
Each $F_t$ is the satisfying-assignment count of a 3-CNF with
$O(n+m+1)$ variables and $O(J+m+1)$ clauses.
The formula can be built in polynomial time.
\end{theorem}
\begin{proof}
Use color variables $x_1,\ldots,x_n$ and one selector $y_i$ per edge.
For each nonempty edge $e_i$, choose an anchor $a_i\in e_i$.
For every $v\in e_i\setminus\{a_i\}$, add
\begin{equation}\label{eq:guards}
 (\neg y_i\vee\neg x_{a_i}\vee x_v),\qquad
 (\neg y_i\vee x_{a_i}\vee\neg x_v).
\end{equation}
These clauses say that selecting an edge forces all its vertices to have
the same color.
Empty and singleton edges need no clauses.
Add the exact-$t$ selector count from Lemma~\ref{lem:adder}.

Fix a coloring $x$.
Its allowed selector assignments are precisely the $t$-element subsets of
its $u(x)$ monochromatic edges.
There are $\binom{u(x)}t$ such assignments, each with one extension to
the gate variables.
Summing over $x$ proves that the count is $F_t$.

The edge constraints use at most $2J$ clauses and no extra variables.
The selector count adds $O(m+1)$ variables and clauses.
The encoding length is
$O((J+n+m+1)\log(n+m+2))$.
\end{proof}

The selectors use one-way implications: a monochromatic edge may be
selected or left unselected.
The gate constraints use equivalences so their extra variables do not
multiply the count.

\section{Running time}

\begin{theorem}[Complete distribution]\label{thm:profile}
Assume $\Count_\alpha$.
All coefficients of $P_H$ can be computed deterministically in polynomial
space and time
\begin{equation}\label{eq:profile}
 (L+1)^{O(1)}2^{O(f_\alpha(n+m))}.
\end{equation}
\end{theorem}
\begin{proof}
Build the formula in Theorem~\ref{thm:encoding} for each $t=0,\ldots,m$
and count its satisfying assignments.
Each formula has $v\le C(n+m+1)$ variables for an absolute constant $C$.
Its length is polynomial in $L+1$, and
$f_\alpha(C(s+1))=O(f_\alpha(s))$.
The $m+1$ calls and Lemma~\ref{lem:inverse} add only polynomial work.
Run the calls one at a time.
Their outputs and all intermediate sums have polynomial bit length,
so the total space is polynomial.
\end{proof}

\begin{corollary}[Set Splitting]\label{cor:parameter}
Under $\Count_\alpha$, Set Splitting with target $k$ can be decided in
polynomial space and time
\[
 |I|^{O(1)}+2^{O(f_\alpha(k))}.
\]
Thus the revised SAT manuscript gives
$2^{O(k\log\log k/\sqrt{\log k})}+|I|^{O(1)}$.
The scope note gives
$2^{O(k(\log\log k)^2/\sqrt{\log k})}+|I|^{O(1)}$.
Both bounds are subexponential in $k$.
\end{corollary}
\begin{proof}
Apply the decision kernel of Lokshtanov and Saurabh~\cite[Theorem~1]{LS09}.
It answers the instance or returns an equivalent instance $(H',k_0)$
with at most $k$ vertices and $2k$ edges, where $k$ is the original target.
Use Theorem~\ref{thm:profile} on $H'$.
Answer yes if a coefficient at an exponent at least $k_0$ is positive.
The reduced input has polynomial length in $k+1$.
Its polynomial factors are absorbed in the displayed exponential term.
\end{proof}

This kernel preserves the decision, not necessarily the number of
colorings. The corollary counts the reduced instance.
Theorem~\ref{thm:profile} counts the original instance in terms of $n+m$.

The SAT manuscript also states a faster bound
$\poly(\ell+v)2^{O_D(v/\sqrt{\log v})}$ for formulas with at most $Dv$
clauses, for fixed $D$.
If $J=O(n+m)$, our formulas have fixed clause density, giving
\[
 (L+1)^{O(1)}2^{O((n+m)/\sqrt{\log(n+m)})}.
\]
This faster rate does not apply to every kernel: it may have
$\Theta(k^2)$ incidences and only $O(k)$ variables.

\section{Extensions}

\paragraph{Fixed colors and an optimal coloring.}
Unit clauses fix vertex colors, so the same method counts their extensions.
Let $b$ be the largest exponent with a positive coefficient.
Fix each vertex to zero if the coefficient at $b$ stays positive; otherwise
fix it to one. An optimum remains possible at every step.
At most $n+1$ computations recover an optimal coloring within
\eqref{eq:profile}, up to a polynomial factor.

\paragraph{Families of 2-CNF formulas.}
Let $\Phi_1,\ldots,\Phi_m$ be 2-CNF formulas on the same $n$ variables,
with $C$ clauses in total.
For every clause $D$ of $\Phi_i$, add $\neg y_i\vee D$.
Selecting $i$ now forces $\Phi_i$ to hold.
The exact-$t$ selector count and binomial inversion recover the number
of assignments making exactly $j$ of the formulas true, for every $j$.
Each SAT call uses $O(n+m+1)$ variables and $O(C+m+1)$ clauses, so
Theorem~\ref{thm:profile} applies.

A clause is false exactly when all its literals are false, a conjunction
of unit clauses. This also gives the satisfaction distribution for any CNF,
regardless of clause width.

\section{Checks and scope}

Run \texttt{python3 research/verify\_paper.py} from the repository root.
The code checks the SAT counts against exhaustive assignments, including
repeated edges, unused vertices, fixed colors, and the extensions above.
Results and source hashes are in \path{paper/verification.json}.

The tests check the reduction, not the external SAT algorithm.
The polynomial identity is established prior work~\cite{White11}.
The open question is whether arbitrary edge sizes admit the faster
fixed-density bound with only $O(n+m)$ variables.

\begin{thebibliography}{9}
\bibitem{LS09}
Daniel Lokshtanov and Saket Saurabh.
\newblock Even Faster Algorithm for Set Splitting!
\newblock In \emph{Parameterized and Exact Computation}, LNCS 5917,
pp.~288--299, Springer, 2009.
\newblock \href{https://doi.org/10.1007/978-3-642-11269-0_24}{doi:10.1007/978-3-642-11269-0\_24}.

\bibitem{NRD14}
Jesper Nederlof, Johan M.~M. van Rooij, and Thomas~C. van Dijk.
\newblock Inclusion/Exclusion Meets Measure and Conquer.
\newblock \emph{Algorithmica}, 69(3):685--740, 2014.
\newblock \href{https://doi.org/10.1007/s00453-013-9759-2}{doi:10.1007/s00453-013-9759-2}.

\bibitem{SAT26}
OpenAI.
\newblock \emph{The Exponential Time Hypothesis is False}.
\newblock Preprint, October 2026; revised October 25, 2026.
\newblock \href{https://github.com/openai/math/blob/main/preprints/The-Exponential-Time-Hypothesis-is-False-October-25-2026/build/main_full_version.tex}{Manuscript}, main theorem. Accessed October 29, 2026.

\bibitem{Scope26}
OpenAI.
\newblock \emph{The Exponential Time Hypothesis is False: Lean formalization scope}.
\newblock 2026.
\newblock \href{https://github.com/openai/math/blob/main/lean/docs/378.md}{Formalization scope note}.
Accessed October 29, 2026.

\bibitem{White11}
Jacob~A. White.
\newblock On Multivariate Chromatic Polynomials of Hypergraphs and Hyperedge Elimination.
\newblock \emph{Electronic Journal of Combinatorics}, 18(1):P160, 2011.
\newblock \href{https://doi.org/10.37236/647}{doi:10.37236/647}.
\end{thebibliography}
\end{document}
